\begin{abstract}
\draft{Draft abstract. Numbers marked TODO come from the Qwen3 math and coding
runs and must be filled from the final experiment logs.}
Reinforcement learning with verifiable rewards is bounded by memory before it
is bounded by compute: on a single 48\,GB GPU, dense-attention GRPO on a 7B
model runs out of memory at 24k tokens of context. Prior work accelerates the
rollout with sparse attention while keeping the policy gradient dense, which
introduces an actor--policy mismatch that must be managed to avoid collapse.
We take the opposite route and make the \emph{entire} RL step sparse: both
generation and the chunked backward pass run through a bounded, actively
evicting KV cache, so the sampled log-probabilities and the gradients come
from the same model and no importance correction is needed. The footprint is
flat in context length (35.7\,GB from 24k to 65k tokens for Qwen2.5-7B), which
lets us train GRPO and RLOO policies at context lengths where dense training
cannot run at all. We train Qwen3 thinking models with verifiable rewards on
Polaris (math) and TACO (code) and evaluate on AIME 2024--2026, AMC, MATH500,
LiveCodeBench, HumanEval+ and MBPP+. Where dense training fits, training
through the evicting cache matches it within noise at lower memory. Past the
dense limit, \draft{TODO: accuracy gained over the base model at a
generation length dense GRPO cannot reach on the same GPU}. We characterise
what eviction destroys and what RL can and cannot recover, and release the
training loop, prompt-set manifests, and cache configurations.
\end{abstract}
